\documentclass[12pt]{article}

% Prevent words to overflow over the margin
\emergencystretch 3em
\usepackage[margin=1.2in]{geometry}

% Headers, packages, definitions
% Images and plots
\usepackage[usenames,dvipsnames]{xcolor}
\usepackage{graphicx, pgfplots, subfigure, standalone, pgf, tikz}
\usepackage[justification=centering]{caption}

%Biblio
\usepackage[sort&compress,numbers]{natbib}

% Utilities
\usepackage{comment, url,float, setspace, xspace}
\pgfplotsset{compat=1.18} 

% Math
\usepackage{amsmath, amssymb, amsthm, amsfonts, enumitem, mathtools, array, bm}
\renewcommand{\qedsymbol}{$\blacksquare$}
\renewcommand\arraystretch{1.2}
\setlist[enumerate]{label=(\roman*.)}

% Algorithms
\usepackage[ruled,vlined,linesnumbered]{algorithm2e}

% Definition of Optimization problems
\usepackage[short,c2,nocomma]{optidef}

% Hyperreferences
\usepackage[colorlinks,
            linkcolor=BlueViolet,
            citecolor=blue,
            urlcolor=blue]{hyperref}

% References to objects
\usepackage[nameinlink]{cleveref}

% Tables
\usepackage{csvsimple, booktabs, longtable, multirow, multicol, adjustbox}	% reading CSV files in tables

% Acronyms
\usepackage[xindy, acronyms, nowarn]{glossaries}   % Glossaries and Acronyms
\makeglossaries

% Todos
\newcommand{\todo}[1]{ {\color{Turquoise} \texttt{\#TODO }#1} }


%ORCIDs
\definecolor{lime}{HTML}{A6CE39}
\DeclareRobustCommand{\orcidicon}{
	\begin{tikzpicture} \draw[lime, fill=lime] (0,0) circle [radius=0.16] node[white] { {\fontfamily{qag}\selectfont \tiny ID} };
	\draw[white, fill=white] (-0.0625,0.095) circle [radius=0.007];
	\end{tikzpicture} \hspace{-2mm}
}

\foreach \x in {A, ..., Z}{\expandafter\xdef\csname orcid\x\endcsname{\noexpand\href{https://orcid.org/\csname orcidauthor\x\endcsname}{\noexpand\orcidicon} } }
% Bibliography
\bibliographystyle{plainnat}

% Definitions
% Sets
\newcommand{\xofi}{x^i}
\newcommand{\xminusi}{x^{-i}}
\newcommand{\X}[1]{ \mathcal{X}^{ #1 } } 
\newcommand{\ones}{\mathbf 1}
\newcommand{\R}{\mbox{\bf R}}
\newcommand{\Z}{\mbox{\bf Z}}
\newcommand{\Sy}{\mbox{\bf S}} 
\newcommand{\NPhard}{\mbox{$\mathcal{NP}$-hard}}  
\newcommand{\SigmaTwoP}{\mbox{$\Sigma^p_2$-hard}} 
\newcommand{\prob}{ {\mathbf P} }
\newcommand{\identity}{I}
\newcommand{\nullspace}{ {\mathcal N} }
\newcommand{\range}{ {\mathcal R} }
\newcommand{\Rank}{\operatorname{\bf Rank}}
\newcommand{\Tr}{\operatorname{\bf Tr}}
\newcommand{\diag}{\operatorname{\bf diag}}
\newcommand{\lambdamax}{\lambda_{\rm max} } 
\newcommand{\lambdamin}{\lambda_{\rm min}}
\newcommand{\Expect}{\operatorname{\bf E{}}}
\newcommand{\Prob}{\operatorname{\bf Prob}}
\newcommand{\dist}{\operatorname{\bf dist{}}}
\newcommand{\sol}{\operatorname{SOL}}
\newcommand{\conv}{\operatorname{conv}}
\newcommand{\intr}{\operatorname{int}}
\newcommand{\relintr}{\operatorname{relint}}
\newcommand{\relint}{\relintr}
\newcommand{\rec}{\operatorname{rec}}
\newcommand{\lin}{\operatorname{lin}}
\newcommand{\bd}{\operatorname{bd}}
\newcommand{\cl}{\operatorname{cl}}
\newcommand{\proj}{\operatorname{proj}}
\newcommand{\sgn}{\operatorname{sgn}}
\newcommand{\cone}{\operatorname{cone}}
\newcommand{\epi}{\operatorname{\bf epi}} % epigraph
\newcommand{\Vol}{\operatorname{\bf vol}}
\newcommand{\dom}{\operatorname{\bf dom}} % domain

\newcommand{\sign}{\operatorname{\bf sign}}
\newcommand{\round}{\operatorname{\bf round}}
\newcommand{\define}{\coloneqq}

% Trees
\newcommand{\tree}{\mathcal{T}}

% Strategy
\newcommand{\strategy}{\mathcal{S}}
\newcommand{\intvars}{\mathcal{I}}
\newcommand{\tightconstraints}{\mathcal{T}}

% Machine learning
\newcommand{\loss}{\mathcal{L}}

% Sections (for citation)
\newcommand{\Sec}{Section\;}

% Create theorems and other environments
\newtheorem{theorem}{Theorem}[section]  % Restart counter every section
\newtheorem{lemma}{Lemma}[section]  % Restart counter every section
\newtheorem{corollary}{Corollary}[theorem]  % Restart counter every theorem
\newtheorem{proposition}{Proposition}[section]
\newtheorem{assumption}{Assumption}[section]
\newtheorem{definition}{Definition}[section]
\newtheorem{example}{Example}[section]
\newtheorem{exercise}{Exercise}[section]
\newtheorem{remark}{Remark}[section]
\newtheorem{problem}{Problem}[section]


% Add Clevref names
\crefname{line}{Line}{Lines}
\crefname{lemma}{Lemma}{Lemmata}
\crefname{theorem}{Theorem}{Theorems}
\crefname{proposition}{Proposition}{Propositions}
\crefname{algorithm}{Algorithm}{Algorithms}
\crefname{equation}{Equation}{Equations}
\crefname{definition}{Definition}{Definition}
\crefname{claim}{Claim}{Claim}
\crefname{corollary}{Corollary}{Corollaries}
\crefname{remark}{Remark}{Remarks}
\crefname{example}{Example}{Examples}
\crefname{figure}{Figure}{Figures}
\crefname{section}{Section}{Sections}
\crefname{table}{Table}{Tables}

% Acronyms
\newacronym{LO}{LO}{Linear Optimization}
\newacronym{QO}{QO}{Quadratic Optimization}
\newacronym{MIQO}{MIQO}{Mixed-Integer Quadratic Optimization}
\newacronym{MIO}{MIP}{Mixed-Integer Programming}
\newacronym{MILO}{MILP}{Mixed-Integer Linear Programming}
\newacronym{MINLO}{MINLP}{Mixed-Integer Nonlinear Programming}
\newacronym{AGT}{AGT}{Algorithmic Game Theory}
\newacronym{POA}{PoA}{Price of Anarchy}
\newacronym{POS}{PoS}{Price of Stability}
\newacronym{PNE}{PNE}{Pure Nash Equilibrium}
\newacronym{MNE}{MNE}{Mixed Nash Equilibrium}
\newacronym{IPG}{IPG}{Integer Programming Game}




\newcommand{\orcidauthorA}{ {{cookiecutter.ORCID}} }
\title{ {{cookiecutter.title}} }
\author{ {{cookiecutter.author}} \orcidA{}}

\begin{document}
\maketitle

\begin{abstract}
	Lorem ipsum dolor sit amet, consectetur adipiscing elit. Nunc cursus mi at pretium congue. Suspendisse enim tortor, scelerisque non luctus quis, malesuada eget mauris. Sed quis scelerisque quam, sed facilisis mauris. Vivamus feugiat tempus urna, a tristique sem gravida id. In fringilla luctus risus quis auctor. Nam lorem ex, condimentum id elit eget, ornare luctus sapien. Phasellus fringilla tincidunt diam, ut blandit diam. Pellentesque habitant morbi tristique senectus et netus et malesuada fames ac turpis egestas. Aenean condimentum imperdiet elit et aliquam. Vivamus scelerisque, lorem in pretium placerat, mi enim molestie sem, eu varius lorem ex quis turpis.
\end{abstract}

% Table of contents
%\newpage \tableofcontents \newpage


\section{Introduction}

Several real-world problems often involve a series of selfish agents optimizing their benefits while mutually affecting their decisions. ~\gls{IPG} offer a flexible modeling paradigm for multi-agent decision-making ~\cite{conforti2014}.

\todo{Here is a todo note regarding \cref{def:IPG}.}


\subsection{Literature Review}

\section{Our method}

\begin{definition}
	An ~\gls{IPG} is a non-cooperative simultaneous game among $n$ players with each player $i=1,2,\dots,n$ solving the optimization problem
	\begin{mini!}
		{\xofi}{f^i(\xofi,\xminusi)  \protect\label{eq:IPG:Obj}}
		{\label{eq:IPG}}{}
		%
		\addConstraint{\xofi}{\in \X{i}, \protect\label{eq:IPG:Constraints}}
	\end{mini!}
	where, for any player $i$, $\X{i}$ is the set $i$'s strategies, and $f^i$ is its utility function.
	\label{def:IPG}
\end{definition}


\todo{Here is another todo note.}


\subsection{Theoretical results}%
\label{sub:theoretical_results}




\begin{theorem}
	This is true.
\end{theorem}
\begin{proof}
	That's why.
\end{proof}

\section{Experimental results}
Here is an example table.

\begin{table}
  \centering
  \caption{Example table.}
  \label{tab:tablelabel}
  \begin{tabular}{ll}
    \toprule
    $A$ & $B$\\
    \midrule
    % Add directly from csv reader
    % the csv file has names colnameA, colnameB
    \csvreader[head to column names, late after line=\\]{./data/table.csv}{
    colnameA=\colA,
    colnameB=\colB
    }{\colA & \colB}
    \bottomrule
  \end{tabular}
\end{table}

Figure~\ref{fig:tikz} shows an example of how to use TikZ.
\begin{figure}
	\centering
  \begin{tikzpicture}
    \begin{axis}[
        samples=100, % you don't need 1000, it only slows things down
        ticks=none,
        xmin = -2.5, xmax = 2.5,
        ymin = -1, ymax = 3,
        axis x line=middle,
        axis y line=middle,
	xlabel={$x$},
        ylabel={$y$},
	x label style={
	  at={(axis cs:2,-0.2)},
	  anchor=west,
	},
	declare function={f(\x)=abs(\x);},
        ]
      \addplot[thick,color=black, mark=none, domain=-1:1, -,shorten >=1pt] {f(x)};
      \draw[dashed] (axis cs:-1,0) node[below=0.5mm] {$-1$} -- (axis cs:-1,2);
      \draw[thick] (axis cs:-1,1) -- (axis cs:-1,2);
      \draw[dashed] (axis cs:1,0) node[below=0.5mm] {$1$} -- (axis cs:1,2);
      \draw[thick] (axis cs:1,1) -- (axis cs:1,2);
      \addplot[mark=*,fill=black] coordinates {(0,0)};
      \addplot[mark=*,fill=black] coordinates {(1,1)};
      \addplot[mark=*,fill=black] coordinates {(-1,1)};
      \draw[->, thick](axis cs:2,2)--(axis cs:2.3,2.2) node[below=0.5mm](q){$\alpha$};
    \end{axis}
  \end{tikzpicture}
  \caption{Example of TikZ figure.}
  \label{fig:tikz}
  \end{figure}

Figure~\ref{fig:pgffigure} shows an example of how to plot with PGFPlots
\begin{figure}
	\begin{tikzpicture}
		\begin{axis}
			[
				axis x line=bottom,
				axis y line=left,
				enlarge y limits={0.05, upper},  % Enlarge y limits by 10%
				enlarge x limits={0.05, upper},  % Enlarge x limits by 10%
				xtick = {0,...,5},
				xmin = 0,
				ymin = 0,
				width=\textwidth,
				height=0.33\textheight,
				xlabel=$x$,
				ylabel=$f(x)$,
				]
				\addplot[thick, color=black, mark=none] table [x=a, y=b, col sep=comma] {./data/plot.csv};
				\legend{My Line $f(x)$}
		\end{axis}
	\end{tikzpicture}
	\caption{My Caption.}
	\label{fig:pgffigure}
\end{figure}


% Bibliography
\bibliography{bibliography}

\end{document}
